\subsection{The extracted functionals and their verification}
\label{sec:p07-dialectica-proof}

All tuples below abbreviate finite lists of typed terms. No product
type or simultaneous-recursion primitive is added to the arrow-only
System T. Every finite type has a closed default term: use zero at
base type, and a constant function to the default of the result type
at an arrow type. Empty witness and challenge tuples require no term.

\begin{hblemma}{thm:p07-t-simultaneous-recursion}{proved=annex:p7.d.1}{Tuple recursion}
For a finite list of types \(\sigma_1,\ldots,\sigma_r\), initial
terms \(a_i:\sigma_i\), and typed step terms
\(F_i(n,x_1,\ldots,x_r):\sigma_i\), System T defines terms
\(X_i(n):\sigma_i\) such that \(\mathsf T_{\mathrm{qf}}\) proves
\[
 X_i(0)=a_i,\qquad
 X_i(Sn)=F_i(n,X_1(n),\ldots,X_r(n)).
\]
Parameters of all finite types are permitted. The term construction
and the proofs of these equations are effective.
\end{hblemma}

\begin{hblemma}{thm:p07-t-matrix-decision}{proved=annex:p7.d.2}{Deciding the matrices}
For every Dialectica matrix of an arithmetic formula there is a
System-T term \(d_A\) of ground type, with the same parameters,
such that the target proves
\[
 d_A=0\quad\longleftrightarrow\quad A.
\]
There are typed conditional terms
\(\operatorname{If}_\sigma(e,u,v)\) taking value \(u\) when
\(e=0\), and \(v\) otherwise. These constructions test only
numerical equalities. They do not decide equality at arbitrary
functional types.
\end{hblemma}

\begin{hblemma}{thm:p07-t-backward-induction}{proved=annex:p7.d.3}{Backward challenges}
Let \(P(n,\vec y)\) be a quantifier-free target formula, with
\(n\) of ground type and a finite typed tuple \(\vec y\).
If \(\mathsf T_{\mathrm{qf}}\) proves
\[
 P(0,\vec y),\qquad
 P(n,\vec b(n,\vec y))\longrightarrow P(Sn,\vec y),
\]
then it proves \(P(n,\vec y)\). The proof transformation uses only
System-T terms, propositional reasoning, and open induction. In
particular it does not instantiate a pointwise induction hypothesis
at a different unlicensed value of its challenge parameter.
\end{hblemma}

\begin{hblemma}{thm:p07-dialectica-induction}{proved=annex:p7.d.4}{The induction schema}
Every instance of the HA induction schema has a Dialectica realizer
in System T, with its matrix proved in
\(\mathsf T_{\mathrm{qf}}\), uniformly in the formula and its
numerical parameters. This treats the full induction axiom of the
handbook's HA presentation, not merely an inference rule for two
already proved closed premises.
\end{hblemma}

The proof of soundness below checks every logical rule, the selected
arithmetic axioms including predecessor existence, and induction.
The conditional for contraction selects a false challenge when one
is available. With the convention \(d_A=0\leftrightarrow A\), it is
\[
 \operatorname{If}(d_{A(\vec x,\vec y_1)},\vec y_2,\vec y_1),
\]
component by component. In the held July 20, 1999 preprint, the first
line of p.~13 prints the last two arguments in the opposite order,
although p.~12 uses the same zero-for-truth convention. That displayed
choice fails when the first challenge makes the matrix true and the
second makes it false. The proof uses the displayed corrected order;
no claim about the wording of the separate 1998 published version is
made from that preprint.
